
\documentstyle[amssymb,preprint,aps]{revtex}
%%%%%%%%%%%%%%%%%%%%%%%%%%%%%%%%%%%%%%%%%%%%%%%%%%%%%%%%%%%%%%%%%%%%%%%%%%%%%%%%%%%%%%%%%%%%%%%%%%%%%%%%%%%%%%%%%%%%%%%%%%%%
%TCIDATA{OutputFilter=LATEX.DLL}
%TCIDATA{Created=Sunday, November 18, 2001 09:49:30}
%TCIDATA{LastRevised=Thursday, November 22, 2001 10:08:53}
%TCIDATA{<META NAME="GraphicsSave" CONTENT="32">}
%TCIDATA{<META NAME="DocumentShell" CONTENT="Articles\SW\REVTeX - APS and AIP Article">}
%TCIDATA{CSTFile=revtxtci.cst}

\newtheorem{theorem}{Theorem}
\newtheorem{acknowledgement}[theorem]{Acknowledgement}
\newtheorem{algorithm}[theorem]{Algorithm}
\newtheorem{axiom}[theorem]{Axiom}
\newtheorem{claim}[theorem]{Claim}
\newtheorem{conclusion}[theorem]{Conclusion}
\newtheorem{condition}[theorem]{Condition}
\newtheorem{conjecture}[theorem]{Conjecture}
\newtheorem{corollary}[theorem]{Corollary}
\newtheorem{criterion}[theorem]{Criterion}
\newtheorem{definition}[theorem]{Definition}
\newtheorem{example}[theorem]{Example}
\newtheorem{exercise}[theorem]{Exercise}
\newtheorem{lemma}[theorem]{Lemma}
\newtheorem{notation}[theorem]{Notation}
\newtheorem{problem}[theorem]{Problem}
\newtheorem{proposition}[theorem]{Proposition}
\newtheorem{remark}[theorem]{Remark}
\newtheorem{solution}[theorem]{Solution}
\newtheorem{summary}[theorem]{Summary}

\begin{document}
\title{A Simple APS RevTeX Article}
\author{J. C. Stenerson}
\address{600 Ericksen Avenue NE, Suite 300, Bainbridge Island, WA 98110}
\date{\today}
\maketitle
\pacs{23.23.+x, 56.65.Dy}

\begin{abstract}
Replace this text with the abstract of your paper.
\end{abstract}

\section{Linear Lagrangians}

We are considering Lagrangians ${\cal L}$ and Lagrangian densities ${\frak L}
$ of the form 
\begin{eqnarray}
{\cal L}\left( \Psi \left( t\right) ,\dot{\Psi}\left( t\right) \right) &=&%
\frac{i}{2}\frac{\left( \left\langle \Psi \left( t\right) |\dot{\Psi}\left(
t\right) \right\rangle -\left\langle \dot{\Psi}\left( t\right) |\Psi \left(
t\right) \right\rangle \right) }{\left\langle \Psi \left( t\right) |\Psi
\left( t\right) \right\rangle }-\frac{\left\langle \Psi \left( t\right)
|H\Psi \left( t\right) \right\rangle }{\left\langle \Psi \left( t\right)
|\Psi \left( t\right) \right\rangle }  \nonumber \\
&=&\frac{i}{2}\frac{\int \left\{ \Psi ^{\ast }\left( {\sf z,}t\right) \dot{%
\Psi}\left( {\sf z,}t\right) -\dot{\Psi}^{\ast }\left( {\sf z,}t\right) \Psi
\left( {\sf z,}t\right) -\Psi ^{\ast }\left( {\sf z,}t\right) \left[ H\Psi %
\right] \left( {\sf z,}t\right) \right\} }{\left\langle \Psi \left( t\right)
|\Psi \left( t\right) \right\rangle }d{\sf z}  \nonumber \\
&\equiv &\int {\frak L}\left( {\sf z,}t\right) d{\sf z}  \label{laden}
\end{eqnarray}%
Parameterizing the wave functions by the real functions ${\cal \varkappa }%
\left( {\sf z},t\right) $ one obtains 
\begin{equation}
{\cal L}\left( \Psi \left( {\cal \varkappa }\left( t\right) \right) ,\dot{%
\Psi}\left( {\cal \varkappa }\left( t\right) \right) \right) ={\cal L}\left( 
{\cal \varkappa }\left( t\right) ,{\cal \dot{\varkappa}}\left( t\right)
\right) =\int {\cal \dot{\varkappa}}\left( {\sf z},t\right) {\cal Z}\left( 
{\sf z},t\right) d{\sf z}-{\cal E}\left( {\cal \varkappa }\left( t\right)
\right)  \label{ladef}
\end{equation}%
where 
\begin{eqnarray}
{\cal Z}\left( {\sf z},t\right) &=&\frac{i}{2}\left\langle \Psi \left(
t\right) |\Psi \left( t\right) \right\rangle ^{-1}\left\{ \left\langle \Psi
\left( t\right) \left| \frac{\delta \Psi \left( t\right) }{\delta {\cal %
\varkappa }\left( {\sf z}\right) }\right. \right\rangle -\left\langle \left. 
\frac{\delta \Psi \left( t\right) }{\delta {\cal \varkappa }\left( {\sf z}%
\right) }\right| \Psi \left( t\right) \right\rangle \right\} \\
{\cal E}\left( \Psi \left( t\right) \right) &=&\frac{{\cal H}\left( \Psi
\left( t\right) \right) }{{\cal N}\left( \Psi \left( t\right) \right) } \\
{\cal H}\left( \Psi \left( t\right) \right) &=&\left\langle \Psi \left(
t\right) |H\Psi \left( t\right) \right\rangle ;\qquad {\cal N}\left( \Psi
\left( t\right) \right) =\left\langle \Psi \left( t\right) |\Psi \left(
t\right) \right\rangle
\end{eqnarray}

The matter Hamiltonian is given by 
\begin{equation}
H=-\frac{1}{2}\tilde{\nabla}^{2}+V
\end{equation}%
where 
\begin{eqnarray}
\tilde{\nabla}^{2} &=&\left( 
\begin{array}{c}
\tilde{\nabla}_{{\bf r}}^{2} \\ 
\tilde{\nabla}_{{\bf R}}^{2}%
\end{array}%
\right) =\tilde{\nabla}\cdot \tilde{\nabla}=\left( 
\begin{array}{cc}
\tilde{\nabla}_{{\bf r}} & \tilde{\nabla}_{{\bf R}}%
\end{array}%
\right) \cdot \left( 
\begin{array}{c}
\tilde{\nabla}_{{\bf r}} \\ 
\tilde{\nabla}_{{\bf R}}%
\end{array}%
\right) \\
\tilde{\nabla}_{{\bf r}} &=&\frac{1}{\sqrt{m_{e}}}\sum_{1\leq j\leq
N_{e}}\nabla _{{\bf r}_{j}}\otimes I_{N_{e}}\left( \check{j}\right) \otimes
I_{N_{n}} \\
\tilde{\nabla}_{{\bf R}} &=&\sum_{1\leq j\leq N_{n}}\frac{1}{\sqrt{M_{j}}}%
\nabla _{{\bf R}_{j}}\otimes I_{N_{n}}\left( \check{j}\right) \otimes
I_{N_{e}} \\
V &=&V_{e}\left( {\bf r}\right) +V_{en}\left( {\bf r,R}\right) +V_{n}\left( 
{\bf R}\right) \\
V_{e}\left( {\bf r}\right) &=&\frac{1}{2}\sum_{1\leq j\neq k\leq N_{e}}\frac{%
1}{\left\| {\bf r}_{j}-{\bf r}_{k}\right\| }\otimes I_{N_{e}}\left( \check{j}%
,\check{k}\right) \otimes I_{N_{n}} \\
V_{en}\left( {\bf r,R}\right) &=&-\sum_{1\leq j\leq N_{e};1\leq k\leq N_{n}}%
\frac{Z_{k}}{\left\| {\bf r}_{j}-{\bf R}_{k}\right\| }I_{N_{e}}\left( \check{%
j}\right) \otimes I_{N_{n}}\left( \check{k}\right) \\
V_{n}\left( {\bf r}\right) &=&\frac{1}{2}\sum_{1\leq j\neq k\leq N_{n}}\frac{%
Z_{j}Z_{k}}{\left\| {\bf R}_{j}-{\bf R}_{k}\right\| }\otimes I_{N_{n}}\left( 
\check{j},\check{k}\right) \otimes I_{N_{e}}
\end{eqnarray}%
and 
\begin{equation}
\left[ H\Psi \right] \left( {\sf z,}t\right) =\left[ \tilde{\nabla}\Psi
^{\ast }\right] \left( {\sf z,}t\right) \cdot \left[ \tilde{\nabla}\Psi %
\right] \left( {\sf z,}t\right) +V\left( {\sf z}\right) \Psi ^{\ast }\left( 
{\sf z,}t\right) \Psi \left( {\sf z,}t\right)
\end{equation}

One can show that the Lagrange equation associated with this Lagrangian
coincides with Schr\"{o}dingers equation. Since Schr\"{o}dingers equation is
first order in time giving an intial state $\Psi \left( t_{0}\right) $ is
sufficient to fix the future evolution of the system. It then follows that
the Lagrangian ${\cal L}$ taken as a function of $\Psi ^{\ast }\left( {\sf z,%
}t\right) $ and $\Psi \left( t\right) $ and their time derrivatives contains
an excess of dynamical variables. In order to find the conjugate momenta and
to pass to the Hamiltonian formalism it is necessary then to eliminate from
eq. $\left( \ref{laden}\right) $ the redundant variables.

Let 
\begin{equation}
\Psi \left( {\sf z,}t\right) =\Psi _{r}\left( {\sf z,}t\right) +i\Psi
_{i}\left( {\sf z,}t\right) 
\end{equation}%
and 
\begin{equation}
\Psi \left( {\sf z,}t\right) =\Psi \left( {\cal \varkappa }\left( {\sf z,}%
t\right) \right) =\Psi \left( {\cal V}\left( {\sf z,}t\right) ,\omega \left( 
{\sf z,}t\right) \right) =e^{{\cal V}\left( {\sf z,}t\right) +i\omega \left( 
{\sf z,}t\right) }\Psi _{ref}\left( {\sf z,}t\right) 
\end{equation}%
then 
\begin{eqnarray}
\Psi _{r}\left( {\sf z,}t\right)  &=&\Psi _{r}\left( {\cal V}\left( {\sf z,}%
t\right) ,\omega \left( {\sf z,}t\right) \right) =e^{{\cal V}\left( {\sf z,}%
t\right) }\cos \omega \left( {\sf z,}t\right) \Psi _{ref}\left( {\sf z,}%
t\right)   \nonumber \\
\Psi _{i}\left( {\sf z,}t\right)  &=&\Psi _{i}\left( {\cal V}\left( {\sf z,}%
t\right) ,\omega \left( {\sf z,}t\right) \right) =e^{{\cal V}\left( {\sf z,}%
t\right) }\sin \omega \left( {\sf z,}t\right) \Psi _{ref}\left( {\sf z,}%
t\right) 
\end{eqnarray}%
The generalized conjugate momentum $\Pi _{r}\left( {\sf z,}t\right) $
associated with $\Psi _{r}\left( {\sf z,}t\right) $ is given by 
\begin{equation}
\Pi _{r}\left( {\sf z,}t\right) =\frac{\delta {\cal L}\left( \Psi \left(
t\right) ,\dot{\Psi}\left( t\right) \right) }{\delta \dot{\Psi}_{r}\left( 
{\sf z}\right) }
\end{equation}%
where 
\begin{equation}
{\cal L}\left( \Psi \left( t\right) ,\dot{\Psi}\left( t\right) \right) =\int
\Pi _{r}\left( {\sf z,}t\right) \dot{\Psi}_{r}\left( {\sf z,}t\right) d{\sf %
z-}{\cal E}\left( \Psi \left( t\right) ,\Psi \left( t\right) \right) 
\end{equation}%
and 
\begin{equation}
\dot{\Psi}_{r}\left( {\sf z,}t\right) =\left( {\cal \dot{V}}\left( {\sf z,}%
t\right) -\dot{\omega}\left( {\sf z,}t\right) \tan \omega \left( {\sf z,}%
t\right) \right) \Psi _{r}\left( {\sf z,}t\right) 
\end{equation}%
It can also be shown that 
\begin{equation}
\Pi _{r}\left( {\sf z,}t\right) =2\Psi _{i}\left( {\sf z,}t\right) 
\end{equation}

\end{document}
